% The two frame structures the host deals in, and the length set the
% flexible-data format allows.
% Millimetres throughout: \framefield draws at y = 0 in millimetres.
\begin{tikzpicture}[font=\sffamily\scriptsize]

% ================================================================ classic
\begin{scope}
  \node[signame] at (-2mm,3mm) {classic, 16 bytes};
  \framefield{0}{16}{fid}{identifier\\{\tiny 4 bytes, with flags in the top bits}}
  \framefield{16}{8}{fstuff}{len\\{\tiny 1}}
  \framefield{24}{12}{fstuff}{padding\\{\tiny 3}}
  \framefield{36}{40}{fdata}{data, up to 8 bytes}
  \node[chlabel, anchor=west] at (80mm,3mm) {the structure that has existed for decades};
\end{scope}

% ================================================================ flexible-data
\begin{scope}[yshift=-20mm]
  \node[signame] at (-2mm,3mm) {flexible-data, 72 bytes};
  \framefield{0}{16}{fid}{identifier\\{\tiny 4 bytes}}
  \framefield{16}{8}{fstuff}{len\\{\tiny 1}}
  \framefield{24}{8}{fcrc}{flags\\{\tiny 1}}
  \framefield{32}{8}{fstuff}{pad\\{\tiny 2}}
  \framefield{40}{36}{fdata}{data, up to 64 bytes}
  \node[chlabel, anchor=west] at (80mm,3mm) {one byte more of header, eight times the payload};
  \draw[busstub] (28mm,-1mm) -- (28mm,-7mm);
  \node[tlabel, anchor=north, text width=44mm, align=center] at (28mm,-7mm)
    {the flags byte carries the bit rate switch and the error state indicator};
\end{scope}

% ================================================================ the length set
\node[chlabel, anchor=west] at (0mm,-38mm)
  {and the payload lengths the format allows, which are not every value};
\foreach \i/\n in {0/0, 1/1, 2/2, 3/3, 4/4, 5/5, 6/6, 7/7, 8/8,
                   9/12, 10/16, 11/20, 12/24, 13/32, 14/48, 15/64}{
  % the unit goes inside the braces: {0*8}mm is parsed as one expression
  \node[regcell, minimum width=8mm, text width=7mm] at ({\i*8mm},-48mm) {\n};
}
\node[chlabel, anchor=north west, text width=60mm] at (0mm,-52mm)
  {Above eight bytes the length jumps: twelve, sixteen, twenty, twenty-four,
   thirty-two, forty-eight, sixty-four. A payload of forty bytes is sent in a
   forty-eight byte frame with the remainder padded, which matters when chapter
   11 decides what to put in a frame.};

\node[umlnote, text width=54mm, anchor=north west] at (64mm,-52mm)
  {The transport unit of seventy-two bytes is what a socket read returns for a
   flexible-data frame, and it is how the host code tells the two formats apart:
   the size of what arrived. That is why the socket option in step 2 exists,
   and why a program that does not set it receives only classic frames from a
   bus carrying both.};
\end{tikzpicture}
